\chapter*{结语：一张领域全景图}

\addcontentsline{toc}{chapter}{结语：一张领域全景图}

回望这 38 个条目，知识架构可以画成一张图：

\begin{center}
\begin{tabular}{ccc}
Wilson 1974 & $\longrightarrow$ & HMC / 涂抹 / 蒸馏（第二章）\\[2pt]
$\downarrow$ & & $\downarrow$\\[2pt]
梯度流纲领（第三章） & & LaMET / 伪分布革命（第四章）\\[2pt]
$\downarrow$ & & $\downarrow$\\
\multicolumn{3}{c}{重整化攻坚战：RI/MOM、辅助场、混合方案（第五章）}\\[-1pt]
\multicolumn{3}{c}{$\downarrow$}\\[-1pt]
\multicolumn{3}{c}{胶子 PDF：流化的与大动量的两线合流（第六章）}\\[2pt]
$\swarrow$ & & $\searrow$\\[2pt]
全球拟合：CT18、NNPDF（第七章） & & 随机量子化的新生：AI 采样（第八章）
\end{tabular}
\end{center}

以三点观察收束本选集。

\textbf{其一}，领域的决定性理论事件在机器上廉价而在思想上昂贵。准 PDF
提案是一页定义；因子化定理是一个 OPE 论证；流的有限性是一个幂计数
证明。它们之所以决定性，是因为每一个都把一个概念上的不可能——"格点
够不到光锥 $x$ 依赖"、"复合算符需要无穷多反项"——转化为有限、可计算的
纲领。

\textbf{其二}，会合胜过竞争。准分布与伪分布、流化算符与大动量算符、
Hessian 与蒙特卡罗拟合：每一对对手最终都是同一对象的两个坐标，而其
相互交叉检验如今是领域主要的系统误差工具。

\textbf{第三}，本语料库并未封闭。其最新的一层——编排上述各章所记录之
确定性流水线的智能体工作流——把这些论文当作可自动化的领域逻辑而非被
取代之物。这一层能否在未来版本中赢得属于自己的一章，取决于本次选编
赖以进行的同样标准：可复现性、误差控制，以及物理优先。
